\subsection{Subespacio generado}

\begin{frame}
    \frametitle{Subespacio generado}

    \begin{definition}
        Sean $v_1,\ldots,v_r$ vectores de un espacio vectorial $V$ sobre $k$.
        El \alert{subespacio generado} por $v_1,\ldots,v_r$ es el conjunto
        de todas sus combinaciones lineales:
        \[
            \langle v_1,\ldots,v_r\rangle
            =\{\lambda_1v_1+\cdots+\lambda_rv_r:\ \lambda_1,\ldots,\lambda_r\in k\}.
        \]
    \end{definition}

    \medskip

    El subespacio generado por el conjunto vacío es el trivial
    $\langle\varnothing\rangle=\{0\}$.

    \medskip

    El plano $W=\{(v_1,v_2,v_3)\in\R^3: v_1-2v_2+v_3=0\}$
    que acabamos de estudiar es
    \[
        W=\langle(2,1,0),(-1,0,1)\rangle.
    \]
\end{frame}

\begin{frame}
    \frametitle{El subespacio generado es un subespacio}

    \begin{proposition}
        $\langle v_1,\ldots,v_r\rangle$ es un subespacio de $V$ que contiene
        a $v_1,\ldots,v_r$.
    \end{proposition}

    \begin{proof}
        Contiene al vector cero, $0=0v_1+\cdots+0v_r$, luego no es vacío.
        Sean
        \[
            u=\lambda_1v_1+\cdots+\lambda_rv_r,
            \qquad
            w=\mu_1v_1+\cdots+\mu_rv_r
        \]
        dos de sus elementos y $\alpha\in k$. Agrupando términos,
        \[
            u+w=(\lambda_1+\mu_1)v_1+\cdots+(\lambda_r+\mu_r)v_r,
            \qquad
            \alpha u=(\alpha\lambda_1)v_1+\cdots+(\alpha\lambda_r)v_r,
        \]
        que vuelven a ser combinaciones lineales de $v_1,\ldots,v_r$.

        Finalmente, cada $v_i$ es la combinación lineal con coeficiente $1$
        en $v_i$ y $0$ en los demás.
    \end{proof}
\end{frame}

\begin{frame}
    \frametitle{El menor subespacio que contiene a unos vectores}

    \begin{proposition}
        Si $W\subseteq V$ es un subespacio y $v_1,\ldots,v_r\in W$, entonces
        \[
            \langle v_1,\ldots,v_r\rangle\subseteq W.
        \]
    \end{proposition}

    \begin{proof}
        Todo elemento de $\langle v_1,\ldots,v_r\rangle$ es una combinación
        lineal de $v_1,\ldots,v_r$, y los subespacios son cerrados para
        combinaciones lineales.
    \end{proof}

    \medskip

    Así, $\langle v_1,\ldots,v_r\rangle$ es el \alert{menor subespacio}
    de $V$ que contiene a $v_1,\ldots,v_r$: los contiene, y está contenido
    en cualquier otro subespacio que los contenga.
\end{frame}

\begin{frame}
    \frametitle{Generadores de un subespacio}

    Como los subespacios son espacios vectoriales, tiene sentido decir que
    unos vectores \alert{generan} un subespacio.

    \begin{proposition}
        \begin{enumerate}
            \item $\langle v_1,\ldots,v_r\rangle$ está generado por
            $v_1,\ldots,v_r$.
            \item Si un subespacio $W\subseteq V$ está generado por
            $v_1,\ldots,v_r$, entonces $W=\langle v_1,\ldots,v_r\rangle$.
        \end{enumerate}
    \end{proposition}

    \begin{proof}
        (1) Los vectores $v_1,\ldots,v_r$ están en
        $\langle v_1,\ldots,v_r\rangle$, y todo elemento de
        $\langle v_1,\ldots,v_r\rangle$ es, por definición, combinación
        lineal de ellos.

        (2) Que $v_1,\ldots,v_r$ generen $W$ significa que están en $W$ y que
        todo vector de $W$ es combinación lineal de ellos, es decir,
        $W\subseteq\langle v_1,\ldots,v_r\rangle$. La otra inclusión se debe
        a que $\langle v_1,\ldots,v_r\rangle$ es el menor subespacio que
        contiene a $v_1,\ldots,v_r$.
    \end{proof}
\end{frame}

\begin{frame}
    \frametitle{Rectas y planos generados}

    \begin{example}
        Sean $u,w\in\R^3$.
        \begin{itemize}
            \item Si $u\neq0$, $\langle u\rangle=\{\lambda u:\lambda\in\R\}$
            es la \alert{recta} que pasa por el origen con dirección $u$.
            \item Si $u$ y $w$ no son proporcionales, $\langle u,w\rangle$
            es el \alert{plano} que pasa por el origen y contiene a $u$ y $w$.
            \item Si $w=\alpha u$, entonces
            $\lambda u+\mu w=(\lambda+\mu\alpha)u$, luego
            \[
                \langle u,w\rangle=\langle u\rangle
            \]
            es una recta, aunque haya dos generadores.
        \end{itemize}
    \end{example}

    \medskip

    El subespacio generado depende de cómo son los vectores, no solo de
    cuántos hay.
\end{frame}

\begin{frame}
    \frametitle{Generadores de un subespacio de polinomios}

    \begin{example}
        Sea $W=\{p(x)\in k_2[x]:\ p(1)=0\}$. Entonces
        \[
            W=\langle x-1,\ x^2-1\rangle.
        \]
    \end{example}

    \begin{proof}[Comprobación]
        \alert{$\supseteq$.} Los polinomios $x-1$ y $x^2-1$ se anulan en $1$,
        luego están en el subespacio $W$, y por tanto
        $\langle x-1,x^2-1\rangle\subseteq W$.

        \alert{$\subseteq$.} Sea $p(x)=a_0+a_1x+a_2x^2\in W$. Como
        $p(1)=a_0+a_1+a_2=0$, se tiene $a_0=-a_1-a_2$, y
        \[
            p(x)=-a_1-a_2+a_1x+a_2x^2=a_1(x-1)+a_2(x^2-1).
        \]
    \end{proof}
\end{frame}

\begin{frame}
    \frametitle{Generadores de un subespacio de matrices}

    \begin{example}
        El subespacio de las matrices simétricas de $M_2(k)$ es
        \[
            \Bigl\langle
            \begin{pmatrix}1&0\\0&0\end{pmatrix},
            \begin{pmatrix}0&0\\0&1\end{pmatrix},
            \begin{pmatrix}0&1\\1&0\end{pmatrix}
            \Bigr\rangle.
        \]
    \end{example}

    \medskip

    Las tres matrices son simétricas. Recíprocamente, una matriz es simétrica
    si y solo si es de la forma
    \[
        \begin{pmatrix}a&b\\b&c\end{pmatrix}
        =a\begin{pmatrix}1&0\\0&0\end{pmatrix}
        +c\begin{pmatrix}0&0\\0&1\end{pmatrix}
        +b\begin{pmatrix}0&1\\1&0\end{pmatrix}.
    \]
\end{frame}

\begin{frame}
    \frametitle{Generadores de las soluciones de un sistema homogéneo}

    \begin{proposition}
        Sea $A\in M_{m\times n}(k)$ y sean $v_1,\ldots,v_s$ las soluciones
        de $Av=0$ obtenidas en el Tema~1 dando el valor $1$ a una incógnita
        libre y $0$ a las demás. Entonces
        \[
            \{v\in k^n:\ Av=0\}=\langle v_1,\ldots,v_s\rangle.
        \]
    \end{proposition}

    \begin{proof}
        En el Tema~1 vimos que toda solución es combinación lineal de
        $v_1,\ldots,v_s$; esto da la inclusión $\subseteq$. La inclusión
        $\supseteq$ se debe a que $v_1,\ldots,v_s$ son soluciones y el
        conjunto de soluciones es un subespacio.
    \end{proof}

    \medskip

    Por tanto, resolver un sistema homogéneo es encontrar un
    \alert{sistema de generadores} de su conjunto de soluciones.
\end{frame}

\begin{frame}
    \frametitle{¿Pertenece un vector al subespacio generado?}

    \begin{proposition}
        Sean $v_1,\ldots,v_r,u\in k^n$ y sea
        $A=(v_1\mid\cdots\mid v_r)\in M_{n\times r}(k)$. Entonces
        \[
            u\in\langle v_1,\ldots,v_r\rangle
            \iff
            \rg(A)=\rg(A\mid u).
        \]
    \end{proposition}

    \begin{proof}
        $u\in\langle v_1,\ldots,v_r\rangle$ significa que existen
        $\lambda_1,\ldots,\lambda_r\in k$ con
        $\lambda_1v_1+\cdots+\lambda_rv_r=u$, es decir, que el sistema
        \[
            A\begin{pmatrix}\lambda_1\\\vdots\\\lambda_r\end{pmatrix}=u
        \]
        es compatible. Por el teorema de Rouché--Frobenius, esto equivale
        a $\rg(A)=\rg(A\mid u)$.
    \end{proof}
\end{frame}

\begin{frame}[allowframebreaks, t]
    \frametitle{Ejemplo: pertenencia al subespacio generado}

    \begin{example}
        ¿Para qué valores de $a\in\R$ se tiene
        $u=(2,3,a,1)\in\langle v_1,v_2\rangle$, con
        \[
            v_1=(1,0,1,2),
            \qquad
            v_2=(0,1,1,-1)?
        \]
    \end{example}

    Escalonamos la matriz ampliada $(v_1\mid v_2\mid u)$ con
    $F_3-F_1-F_2$ y $F_4-2F_1+F_2$:
    \[
        \left(\begin{array}{cc|c}
            1&0&2\\
            0&1&3\\
            1&1&a\\
            2&-1&1
        \end{array}\right)
        \longrightarrow
        \left(\begin{array}{cc|c}
            1&0&2\\
            0&1&3\\
            0&0&a-5\\
            0&0&0
        \end{array}\right).
    \]

    \framebreak

    La matriz de coeficientes tiene rango $2$. La ampliada tiene rango $2$
    si $a=5$ y rango $3$ si $a\neq5$. Por tanto,
    \[
        u\in\langle v_1,v_2\rangle
        \iff
        a=5.
    \]

    \medskip

    En ese caso, la forma escalonada da directamente los coeficientes:
    $\lambda_1=2$, $\lambda_2=3$, es decir,
    \[
        (2,3,5,1)=2v_1+3v_2.
    \]
\end{frame}

\begin{frame}
    \frametitle{Comparar subespacios generados}

    \begin{proposition}
        Sean $v_1,\ldots,v_r,w_1,\ldots,w_s\in V$. Entonces
        \[
            \langle v_1,\ldots,v_r\rangle\subseteq\langle w_1,\ldots,w_s\rangle
            \iff
            v_1,\ldots,v_r\in\langle w_1,\ldots,w_s\rangle.
        \]
    \end{proposition}

    \begin{proof}
        \alert{$\Rightarrow$.} Cada $v_i$ está en
        $\langle v_1,\ldots,v_r\rangle$, luego en
        $\langle w_1,\ldots,w_s\rangle$.

        \alert{$\Leftarrow$.} $\langle w_1,\ldots,w_s\rangle$ es un subespacio
        que contiene a $v_1,\ldots,v_r$, luego contiene al menor subespacio
        que los contiene, que es $\langle v_1,\ldots,v_r\rangle$.
    \end{proof}

    \medskip

    Para probar que dos subespacios generados son \alert{iguales}, basta
    ver que cada generador de uno pertenece al otro.
\end{frame}

\begin{frame}
    \frametitle{Ejemplo: dos sistemas de generadores del mismo subespacio}

    \begin{example}
        En $\R^3$,
        \[
            \langle(1,1,0),(0,1,1)\rangle=\langle(1,2,1),(1,0,-1)\rangle.
        \]
    \end{example}

    \medskip

    Llamemos $v_1=(1,1,0)$, $v_2=(0,1,1)$, $w_1=(1,2,1)$, $w_2=(1,0,-1)$.

    \medskip

    \alert{$\supseteq$.} $w_1=v_1+v_2$ y $w_2=v_1-v_2$ están en
    $\langle v_1,v_2\rangle$.

    \medskip

    \alert{$\subseteq$.} Despejando,
    \[
        v_1=\tfrac12w_1+\tfrac12w_2,
        \qquad
        v_2=\tfrac12w_1-\tfrac12w_2
    \]
    están en $\langle w_1,w_2\rangle$.

    \medskip

    Un mismo subespacio tiene, en general, \alert{muchos} sistemas de
    generadores distintos.
\end{frame}

\begin{frame}
    \frametitle{Generadores redundantes}

    \begin{corollary}
        Sean $v_1,\ldots,v_r,u\in V$. Entonces
        \[
            \langle v_1,\ldots,v_r,u\rangle=\langle v_1,\ldots,v_r\rangle
            \iff
            u\in\langle v_1,\ldots,v_r\rangle.
        \]
    \end{corollary}
\end{frame}
